\chapterOrPart{About}\label{ch:about}
This appendix contains a very short description of \LaTeX\ and explains the content and use of this template.

\section{About \LaTeX}\label{sec:aboutlatex}
\LaTeX\ is software to make high-quality technical and scientific documents. It is not a \enquote{what you see is what you get} (WYSIWYG) editor such as Microsoft Word or OpenOffice. It is a typesetting engine, which means that it works more like a programming language.

\LaTeX\ requires an input file that contains the source code. The structure of that file is shown in \cref{list:latex} and is discussed in more detail \href{https://en.wikibooks.org/wiki/LaTeX/Document_Structure}{here}.

\begin{lstlisting}[language=latexcode,style=latexstyle,label={list:latex},caption={Structure of a LaTeX file}]
 \documentclass*\oarg{options}\marg{class name}*
    % This part is the PREAMBLE, here:
    %  packages are loaded that enhance the capabilities of LaTeX
    %  settings are defined
 \begin{document}
    % This part is the DOCUMENT, here:
    %  content is provided in plain text format
 \end{document}
\end{lstlisting}

\subsection{Document class}
The input file starts with the command \cs{documentclass}\oarg{options}\marg{class name}, where \meta{options} is a placeholder for a list of options, such as for instance font size, paper size, etc. that can be given to customize the behaviour of the document class and \meta{class name} is a placeholder for the name of a class file, recognised by the extension \texttt{.cls}. In a class file the layout standard is defined which tells \LaTeX\ how to format your content. 

This template uses the class file \pkg{memoir} but many other class files exist such as \pkg{article}, \pkg{report}, \pkg{book},\dots. A comprehensive list of class files can be consulted \href{https://ctan.org/topic/class}{at CTAN}.

\subsection{Preamble}
The next part in the input file is called the \textit{preamble}. In this part, packages are loaded that enhance the capabilities of \LaTeX. This is typically done in so-called style files with the extension \texttt{.sty}. The style files that are loaded in this template are found in \texttt{styles/myStyleFiles}. There are three style files that load packages:
\begin{itemize}
    \item \texttt{styles/myBibSettings}: to load the \pkg{biblatex} package to specify settings related to citing and the bibliography
    \item \texttt{styles/myListings}: to load the  \pkg{listings} package to specify settings related to formatting computer code
    \item \texttt{styles/myPackages}: to specify settings for other packages
\end{itemize}

The purpose and use of some of the loaded packages are discussed further in this chapter. Currently there exist around 4000 packages for \LaTeX\ for all kind of applications. The comprehensive list of all packages can be consulted \href{https://www.ctan.org/pkg/}{at CTAN}.

In the \textit{preamble} a user can also define his/her own commands and environments. In this template, this is done in the other style files in the folder \texttt{styles}.

\subsection{Entering content}
The actual content of your work is enclosed between the commands \cs{begin\{document\}} and \cs{end\{document\}}. When you have a lot of content, you will typically divide it over several files. Files with content have the extension \texttt{.tex}. Content is added as plain text, combined with embedded \LaTeX\ commands for specific typeset results. To include files in your input file, you can use the command: \cs{include\{\meta{filename}\}}  or \cs{input\{\meta{filename}\}}.

The command \cs{include\{\meta{filename}\}} is used to include chapters. It will first process all the floats (figures and tables) from the previous chapter, cause a page break so that the next chapter can start on a new page and import the contents from the file named \meta{filename} into the target file. It will also start a new auxiliary file in the background to keep track of the locations of figures, tables and equations within the chapter for proper cross-referencing. Inspect the file \texttt{mainReport} to see what chapters are included in this template. 

The command \cs{input\{\meta{filename}\}} only imports the contents from the file named \meta{filename} into the target file. It's equivalent to typing all the content from the file \meta{filename} into the target file at the location of the \cs{input\{\meta{filename}\}} command.

The content is typically structured for the reader by using commands such as:\\
\cs{chapter\{\meta{name of chapter}\}}\\
\cs{section\{\meta{name of section}\}}\\
\cs{subsection\{\meta{name of subsection}\}}\\
\cs{subsubsection\{\meta{name of subsubsection}\}}\\

More info on how to add content such as text, images, tables and math can be found in \href{https://www.overleaf.com/learn/latex/Learn_LaTeX_in_30_minutes}{Learn \LaTeX\ in 30~minutes}.

\subsection{Compiling your input file}
You have to compile you input file to obtain a typeset PDF document. The processing of your input file is done by a so-called TeX engine. It will parse your input file, line by line, detect the embedded \LaTeX\ commands in your input file and act upon them, according the settings in your class file and packages. This means that you only have to concentrate on the content, while \LaTeX\ will take care of the formatting.

There are four possible outcomes of a compilation:
\begin{description}
    \item [Successful compilation] This is the desired outcome. The compilation process completed without warning or errors. The output PDF corresponds to the specified content.
    \item [Compilation with warnings] The compilation process completed but there were minor issues in your code. The feedback button in Overleaf will colour orange and show the number of warnings. If there are 3 warnings for instance, you will see \colorbox{orange}{\textcolor{white}{3}}. Click on the box to get more info about the warning. Warnings can be caused by duplicate labels, text extending beyond margins, etc. The output PDF might have issues with formatting, cross-referencing or citing. You should review your code and address the issues that gave rise to the warnings. This might not be an immediate necessity, but it's certainly important to do so before finalizing your document.
    \item [Compilation with errors] The compilation was halted because serious issues were encountered that could not be resolved. The feedback button in Overleaf will colour red and show the number of errors. If there are 8 errors for instance, you will see \colorbox{red}{\textcolor{white}{8}}. Click on the box to get more info about the errors. Errors can be caused by syntax errors, missing files, incorrect commands, etc. You must immediately address these errors because the produced output PDF can not be trusted or will even be missing. Good to know: an error can induce a cascade of other errors. So don't be alarmed by the potential abundance of reported errors. Always start by addressing the first listed error. Once rectified, you might find that other errors also vanish. See the \href{https://www.overleaf.com/learn/latex/Errors}{help pages} of Overleaf for more info on compilation errors.
    \item [Compile time out warning]\label{sec:timeout} In 2023, Overleaf reduced the compile time for free plan users from 1 minute to 20 seconds. That is not enough to compile this entire template. Possible workarounds are:
    \begin{itemize}
        \item Comment out some of the large appendices in the template and compile again. This will reduce the required compilation time.
        \item Or compile the file \texttt{myThesis} instead. This file contains the full structure of a thesis but without the appendices, so it takes less time to compile.
        \item Transfer the ownership of your project to your supervisor or promoter. If he/she has a premium account on Overleaf, enough compilation time will be available to compile the entire template.
        \item Download the template and compile it offline.
        \item Get a paid Overleaf subscription.
    \end{itemize}

\end{description}

\section{About this template}
\subsection{Content of the template}
This template will get you started with writing your thesis, report or slides in \href{https://www.latex-project.org}{\LaTeX}. The template contains:
\begin{itemshort}
    \item a title page for VUB, BRUFACE and ULB
    \item a project plan
    \item a jury (for a PhD)
    \item acknowledgements
    \item an abstract in English, Dutch and French
    \item a disclosure on the use of AI
    \item a table of contents
    \item a list of figures, tables and listings
    \item a nomenclature
    \item the typical chapters
    \item appendices with useful information on various topics
    \item your own publications (for a PhD)
    \item a bibliography
    \item a demo bibliography with common entry types
    \item a to-do list
\end{itemshort}

\subsection{Content of the appendices}
In the appendices you will find useful information about
\begin{itemshort}
    \item \LaTeX\ and this template
    \item formatting of figures, tables, equations and source code
    \item useful \LaTeX\ packages and commands
    \item examples of different types of figures and plots
    \item mathematical and Greek symbols
    \item example of a table in appendix
    \item example of an included PDF file
\end{itemshort}

The template uses the class \href{https://ctan.org/pkg/memoir?lang=en}{memoir}. It is a class for typesetting mathematical works as books, reports and articles.

\subsection{How to start using it}
To use this template: surf to \href{https://www.overleaf.com/latex/templates/tips-for-writing-in-latex/kxtvhpxbmdrs}{this template on the Overleaf website} and click on the button 'Open as Template'. This will create a project with the name \textit{Tips for writing (in LaTeX)}. You can then change the name of the project according to your choice. FYI: the name \textit{Tips for writing (in LaTeX)} was inherited from the precursor to this template, a document with tips on formatting numbers, units, and figures, hence the name.

A \href{https://vub.cloud.panopto.eu/Panopto/Pages/Sessions/List.aspx?folderID=eb5cbed1-a7f8-4d08-bca3-ad1e00db8d7a}{set of videos} is available that demonstrate the use of this template.

This template can generate a report and slides from the same source code. It comes with three main files:
\begin{itemize}
    \item \texttt{mainReport}: to compile a report and all the appendices of this template
    \item \texttt{mainSlides}: to compile slides. Only selected content is put on the slides though. To see all content, you must compile \texttt{mainReport} and/or inspect the source code.
    \item \texttt{PhD}: to compile a PhD but without the appendices of this template
    \item \texttt{Thesis}: to compile a thesis but without the appendices of this template
\end{itemize}

\subsection{How to change the language}
This template is written in English. If you write in another language, you have to adapt the language in two ways (example below is for Dutch):
\begin{enumerate}
    \item Select \texttt{Dutch} as spell check language in Overleaf, via the \texttt{Menu} in the upper left corner. Words that do not belong to the selected spell check language will be underlined in red in the editor. The selected spell checking language has no effect on the output PDF. 
    \item Replace the option \texttt{english} by \texttt{dutch} in the command \cs{documentclass}. Options supplied to \cs{documentclass} are global options and are available to all packages. 
\end{enumerate}

The option \texttt{dutch} will be transferred to the following two packages, loaded in \texttt{styles/myPackages}:
    \begin{itemize}
        \item \verb|\usepackage{babel}|\\
        This will automatically translate standard names:\\
        chapter $\to$ hoofdstuk\\
        contents $\to$ inhoudsopgave\\
        bibliography $\to$ bibliografie\\\dots
        \item \verb|\usepackage{cleveref}|\\
        This will automatically translate cross-referencing names:
        \\figure $\to$ figuur\\
        table $\to$ tabel\\
        equation $\to$ vergelijking\\\dots
    \end{itemize}

The option \texttt{dutch} will also cause the nomenclature to be displayed in Dutch. If you want to know how this works, inspect the first lines of the file \texttt{extra/nomenclature}.

% more info on \mode<all> in section 21.3 of beameruserfile.pdf 
\mode<all>
